\documentclass[11pt,a4paper]{article}

%========================= FONTES E CODIFICAÇÃO =========================
\usepackage[T1]{fontenc}
\usepackage[utf8]{inputenc}
\usepackage{mathptmx}
\usepackage[scaled=0.88]{helvet}
\usepackage{courier}
\usepackage{microtype}

% Sem padrões de hifenização em português nesta instalação:
% penaliza hifenização para evitar quebras incorretas.
\hyphenpenalty=9000
\exhyphenpenalty=9000
\tolerance=2500
\emergencystretch=3em

%========================= GEOMETRIA =========================
\usepackage[a4paper,top=2.6cm,bottom=2.6cm,left=2.7cm,right=2.7cm,headheight=14pt]{geometry}

%========================= PACOTES =========================
\usepackage{xcolor}
\usepackage{graphicx}
\usepackage{booktabs}
\usepackage{tabularx}
\usepackage{longtable}
\usepackage{array}
\usepackage{multirow}
\usepackage{enumitem}
\usepackage{titlesec}
\usepackage{fancyhdr}
\usepackage{caption}
\usepackage{siunitx}
\usepackage[most]{tcolorbox}
\usepackage{amsmath,amssymb}
\usepackage[hidelinks,pdfusetitle]{hyperref}

%========================= CORES =========================
\definecolor{lescblue}{HTML}{1F3A5F}
\definecolor{lescaccent}{HTML}{2E6DA4}
\definecolor{lescgray}{HTML}{4A4A4A}
\definecolor{lesclight}{HTML}{EEF3F8}
\definecolor{warncolor}{HTML}{8A3324}
\definecolor{warnlight}{HTML}{FBEEEA}
\definecolor{okcolor}{HTML}{2F5D3A}
\definecolor{oklight}{HTML}{EDF4EE}
\definecolor{rulegray}{HTML}{B8C4D0}

%========================= SEÇÕES =========================
\titleformat{\section}
  {\normalfont\Large\bfseries\color{lescblue}}
  {\thesection}{0.7em}{}[{\vspace{2pt}\color{rulegray}\titlerule[0.8pt]}]
\titleformat{\subsection}
  {\normalfont\large\bfseries\color{lescaccent}}
  {\thesubsection}{0.6em}{}
\titleformat{\subsubsection}
  {\normalfont\normalsize\bfseries\color{lescgray}}
  {\thesubsubsection}{0.5em}{}
\titlespacing*{\section}{0pt}{18pt}{9pt}
\titlespacing*{\subsection}{0pt}{13pt}{5pt}
\titlespacing*{\subsubsection}{0pt}{10pt}{4pt}

%========================= CABEÇALHO =========================
\pagestyle{fancy}
\fancyhf{}
\renewcommand{\headrulewidth}{0.4pt}
\renewcommand{\footrulewidth}{0pt}
\fancyhead[L]{\footnotesize\color{lescgray}Port do sensor de \emph{aging} para Intel MAX 10}
\fancyhead[R]{\footnotesize\color{lescgray}LESC/UFC}
\fancyfoot[C]{\footnotesize\color{lescgray}\thepage}

%========================= LISTAS =========================
\setlist[itemize]{leftmargin=1.3em,itemsep=2pt,topsep=4pt,parsep=0pt}
\setlist[enumerate]{leftmargin=1.6em,itemsep=3pt,topsep=4pt,parsep=0pt}

%========================= CAIXAS =========================
\newtcolorbox{keybox}[1][]{
  enhanced, breakable,
  colback=lesclight, colframe=lescaccent,
  boxrule=0.4pt, left=8pt, right=8pt, top=7pt, bottom=7pt,
  arc=1.5pt,
  fonttitle=\bfseries\small\color{white},
  coltitle=white, colbacktitle=lescaccent,
  attach boxed title to top left={xshift=6pt,yshift=-3pt},
  boxed title style={arc=1pt,boxrule=0pt},
  #1}

\newtcolorbox{warnbox}[1][]{
  enhanced, breakable,
  colback=warnlight, colframe=warncolor,
  boxrule=0.4pt, left=8pt, right=8pt, top=7pt, bottom=7pt,
  arc=1.5pt,
  fonttitle=\bfseries\small\color{white},
  coltitle=white, colbacktitle=warncolor,
  attach boxed title to top left={xshift=6pt,yshift=-3pt},
  boxed title style={arc=1pt,boxrule=0pt},
  #1}

\newtcolorbox{okbox}[1][]{
  enhanced, breakable,
  colback=oklight, colframe=okcolor,
  boxrule=0.4pt, left=8pt, right=8pt, top=7pt, bottom=7pt,
  arc=1.5pt,
  fonttitle=\bfseries\small\color{white},
  coltitle=white, colbacktitle=okcolor,
  attach boxed title to top left={xshift=6pt,yshift=-3pt},
  boxed title style={arc=1pt,boxrule=0pt},
  #1}

%========================= AMBIENTE DE FRENTE =========================
\newcommand{\frenteheader}[3]{%
  \vspace{6pt}
  \noindent\begin{tcolorbox}[enhanced,colback=lescblue,colframe=lescblue,
    boxrule=0pt,left=9pt,right=9pt,top=6pt,bottom=6pt,arc=2pt]
    {\color{white}\bfseries\large #1}\\[2pt]
    {\color{white!85!lescblue}\small #2}\\[1pt]
    {\color{white!75!lescblue}\footnotesize\itshape #3}
  \end{tcolorbox}
  \vspace{4pt}}

\newcommand{\campo}[1]{\vspace{6pt}\noindent{\bfseries\color{lescgray}#1}\par\vspace{2pt}}

%========================= TABELAS =========================
\renewcommand{\arraystretch}{1.25}
\newcolumntype{L}[1]{>{\raggedright\arraybackslash}p{#1}}
\newcolumntype{C}[1]{>{\centering\arraybackslash}p{#1}}
\captionsetup[table]{font=small,labelfont={bf,color=lescblue},skip=5pt}
\captionsetup[figure]{font=small,labelfont={bf,color=lescblue}}

\newcommand{\Tabela}{Tabela}
\renewcommand{\tablename}{Tabela}
\renewcommand{\figurename}{Figura}
\renewcommand{\contentsname}{Sumário}

\begin{document}

%========================= CAPA =========================
\begin{titlepage}
\thispagestyle{empty}
\vspace*{1.2cm}

\noindent{\color{rulegray}\rule{\textwidth}{1.2pt}}
\vspace{0.9cm}

\noindent{\Huge\bfseries\color{lescblue}Portabilidade do Sensor de \emph{Aging} Auto-Tuning para FPGA Intel MAX 10}

\vspace{0.7cm}
\noindent{\LARGE\color{lescaccent}Plano de Fase de Estudo e Divisão de Tarefas}

\vspace{0.5cm}
\noindent{\color{rulegray}\rule{\textwidth}{1.2pt}}

\vspace{1.4cm}

\noindent\begin{tcolorbox}[enhanced,colback=lesclight,colframe=lescaccent,boxrule=0.4pt,
  left=12pt,right=12pt,top=10pt,bottom=10pt,arc=2pt]
\small
\textbf{Contexto.} Migração da plataforma de sensoriamento de envelhecimento por
\emph{slack}, validada em Artix-7 (XC7A100T, Nexys~4~DDR) e publicada em SBCCI 2025
e SBCCI 2026, para o dispositivo Intel MAX~10 (10M50DAF484) com fluxo Intel Quartus
Prime Lite.

\vspace{5pt}
\textbf{Destinatário.} Equipe de alunos de graduação e pós-graduação sob orientação,
com alocação prevista de 4 a 6 integrantes ao longo de um semestre letivo.

\vspace{5pt}
\textbf{Natureza do documento.} Plano de execução. Define a fase de estudo obrigatória,
a decomposição do trabalho em frentes, os critérios de aceite de cada frente, o
sequenciamento com portões de decisão e a política de risco.
\end{tcolorbox}

\vfill

\noindent\begin{tabularx}{\textwidth}{@{}lX@{}}
\toprule
\textbf{Laboratório} & LESC --- Laboratório de Engenharia de Sistemas de Computação \\
\textbf{Instituição} & Universidade Federal do Ceará (UFC) \\
\textbf{Programa} & PPGETI --- Engenharia Elétrica e de Computação \\
\textbf{Versão} & 1.0 \\
\textbf{Horizonte} & 16 semanas (um semestre letivo) \\
\bottomrule
\end{tabularx}

\vspace{0.8cm}
\noindent{\color{rulegray}\rule{\textwidth}{0.8pt}}
\end{titlepage}

%========================= SUMÁRIO =========================
\pagenumbering{roman}
\setcounter{page}{1}
\tableofcontents
\clearpage
\pagenumbering{arabic}
\setcounter{page}{1}

%=========================================================================
\section{Enquadramento}
%=========================================================================

\subsection{Objetivo}

Portar a plataforma de sensoriamento de envelhecimento por \emph{slack} --- hoje
implementada e validada em FPGA Xilinx Artix-7 (XC7A100T, placa Nexys~4~DDR) ---
para o dispositivo Intel MAX~10 (10M50DAF484, placa DE10-Lite ou equivalente),
utilizando o fluxo Intel Quartus Prime Lite.

O objetivo explícito da migração não é reproduzir os resultados numéricos já
publicados. É produzir uma segunda implementação, em nó tecnológico e fabricante
distintos, que satisfaça quatro requisitos:

\begin{enumerate}
  \item medir o \emph{slack} de um caminho crítico sintético por meio de uma
        referência temporal rastreável e caracterizada;
  \item preservar operacional o laço de auto-\emph{tuning}, isto é, o rastreamento
        automático da borda de falha à medida que ela se desloca por temperatura,
        tensão ou degradação permanente;
  \item registrar a fidelidade de medição efetivamente atingida, expressa em
        picossegundos, com incerteza declarada e método de estimação auditável;
  \item permitir afirmar --- ou refutar, com dados --- que a arquitetura do sensor
        é independente de fabricante e de nó tecnológico.
\end{enumerate}

\subsection{O que este projeto não é}

Três delimitações precisam ser comunicadas à equipe na primeira reunião, porque
cada uma corresponde a um modo de falha recorrente em projetos de migração conduzidos
por equipes em formação.

\begin{itemize}
  \item \textbf{Não é uma tarefa de tradução de código.} A porção de HDL que
        precisa ser reescrita é pequena. O trabalho intelectual está em reconstruir
        a \emph{referência de tempo} sobre um recurso de \emph{hardware} com
        propriedades quantitativamente diferentes, e em provar o que se perdeu.
  \item \textbf{Não é uma replicação.} As duas implementações não produzirão
        números comparáveis em valor absoluto. Serão comparáveis em metodologia
        e em métricas normalizadas, e é isso que o plano protege.
  \item \textbf{Não é um projeto cujo sucesso se mede por ``funcionou''.}
        Um sensor que compila, roda e emite números sem caracterização de régua
        temporal não constitui resultado. O objeto científico aqui é a incerteza
        da medida.
\end{itemize}

\subsection{Premissas sobre a versão de referência}

Este plano assume que a versão mais recente do projeto, sobre Artix-7 / Nexys~4~DDR,
é composta pelos subsistemas da Tabela~\ref{tab:subsistemas}. Divergências devem ser
corrigidas antes da distribuição do documento à equipe, porque a alocação de tarefas
deriva diretamente desta decomposição.

\begin{table}[htbp]
\centering
\caption{Decomposição da versão de referência em Artix-7 e destino de cada subsistema no port.}
\label{tab:subsistemas}
\small
\begin{tabularx}{\textwidth}{@{}L{3.3cm} X L{2.6cm}@{}}
\toprule
\textbf{Subsistema} & \textbf{Função na versão de referência} & \textbf{Frente} \\
\midrule
Sensor de \emph{slack} (RTL) &
Caminho crítico sintético amostrado por par de \emph{clocks} com defasagem
controlada; conversão de evento de falha em medida temporal. & F1 \\
Régua temporal &
Deslocamento incremental de fase do MMCM (passo fino $=$ período do VCO $/\,56$),
travado no cristal da placa. & F1 \\
Laço de auto-\emph{tuning} &
Máquina de estados que localiza a borda de falha e a rastreia continuamente,
com histerese e política de reaquisição. & F1 \\
Restrição física &
Preservação de posicionamento e de estrutura lógica contra otimizações da
ferramenta, garantindo comparabilidade entre compilações. & F2 \\
Telemetria PVT &
XADC fornecendo temperatura de junção e $V_\text{CCINT}$ \emph{on-die};
alimenta o \emph{trim} de tensão em malha fechada. & F3 \\
Comunicação &
Fluxo UART do DUT para o \emph{host}, com protocolo de pacote e política
de \emph{timestamp}. & F4 \\
Bancada térmica &
Forno resistivo acionado por SSR, controlador PID sintonizado por
identificação FOPDT e regras SIMC. & F5 \\
Bancada elétrica &
Fonte programável ITECH IT6502D comandada por VISA, com malha de compensação
de tensão realimentada pelo XADC. & F5 \\
Aquisição &
Aplicação Python com registro CSV temporizado, gráficos ao vivo e recuperação
de falhas de USB em campanhas de vários dias. & F4 \\
Análise &
Decomposição de variância PVT do sinal de \emph{slack}, ajuste de leis de
potência e de Arrhenius, caracterização informacional. & F6 \\
\bottomrule
\end{tabularx}
\end{table}

\subsection{Critério de sucesso do projeto}

\begin{okbox}[title={Definição de pronto}]
O projeto é considerado bem-sucedido quando existir, sobre MAX~10, uma medida de
$\sigma_\text{resid}$ expressa em picossegundos, obtida com metodologia idêntica à
dos artigos anteriores, acompanhada de (i) caracterização da régua temporal empregada,
(ii) barra de incerteza, e (iii) declaração explícita do menor degrau de \emph{slack}
resolvível sob critério $k\sigma$.

Este critério é satisfeito \emph{tanto} pelo desfecho favorável quanto pelo desfavorável.
Um resultado que estabeleça ``a arquitetura não porta com fidelidade útil para este nó
tecnológico, pelas razões X e Y, com penalidade quantificada em Z'' encerra o projeto
com sucesso.
\end{okbox}

%=========================================================================
\section{Diagnóstico técnico de portabilidade}
%=========================================================================

Esta seção existe para que a equipe entre na Fase~0 sabendo onde está a dificuldade.
É material de leitura obrigatória, não um resumo executivo.

\subsection{A régua temporal é a origem da resolução}

A observação central, e que reordena toda a hierarquia de dificuldades do port, é
esta: a resolução do sensor na versão Artix-7 \emph{não} vem da malha lógica nem da
linha de atraso, e sim do passo de deslocamento de fase do MMCM.

No MMCM da família 7, o passo fino de deslocamento corresponde ao período do VCO
dividido por 56. Na configuração empregada, isso resulta em aproximadamente
\SI{19.8}{\pico\second} por incremento --- coerente com a relação reportada nos
artigos, em que \SI{0.90}{contagens} correspondem a cerca de \SI{17.9}{\pico\second}.

Essa escolha arquitetural tem uma consequência que costuma passar despercebida e que
é, na verdade, o que torna o sensor válido como instrumento de envelhecimento:
\textbf{a régua está travada no cristal da placa}. O período do VCO é derivado de um
oscilador de quartzo, cuja deriva com temperatura e tensão é ordens de grandeza menor
que a do silício sob medição. O caminho medido envelhece; a régua, não. Um sensor
construído sobre linha de atraso em tecido lógico não tem essa propriedade --- a régua
envelhece junto com o mensurando, e a separação entre os dois exige calibração adicional
em cada ponto de operação.

\begin{keybox}[title={Consequência para o desenho da equipe}]
Qualquer proposta de recuperar resolução no MAX~10 precisa ser avaliada primeiro sob
o critério de \emph{rastreabilidade}, e só depois sob o critério de finura. Uma régua
fina que deriva com PVT é, para o propósito deste projeto, pior do que uma régua grossa
e estável --- a menos que a deriva seja caracterizada e corrigida ponto a ponto,
o que é um projeto em si.
\end{keybox}

\subsection{A restrição dominante do MAX 10}

Os PLLs do MAX~10 suportam reconfiguração dinâmica e deslocamento de fase em tempo de
execução, o que preserva a arquitetura do sensor. A limitação está na granularidade.

Conforme o \emph{MAX~10 Clocking and PLL User Guide}, o menor deslocamento de fase
corresponde ao período do VCO dividido por oito, e a frequência do VCO deve permanecer
entre \SI{600}{\mega\hertz} e \SI{1300}{\mega\hertz} para operação correta do PLL. A
documentação especifica a resolução de fase como incrementos ``de até
\SI{96}{\pico\second}''. Com o VCO no teto da faixa:

\begin{equation}
\Delta t_\text{min} = \frac{1}{8}\cdot\frac{1}{\SI{1300}{\mega\hertz}}
= \frac{\SI{769.2}{\pico\second}}{8} \approx \SI{96.2}{\pico\second}.
\end{equation}

Este é um piso absoluto: não há configuração de contadores, modo de compensação ou
opção de ferramenta que o reduza. Qualquer resolução mais fina exige uma arquitetura
de medição diferente, não uma sintonia diferente.

\subsection{Propagação da perda de resolução}

A penalidade não é uma abstração. Ela se propaga diretamente para a métrica que
sustenta a aplicabilidade do sensor em gerenciamento de ciclo de vida do silício (SLM).
Aplicando o mesmo $\sigma_\text{resid}$ em contagens já medido nos dois bancos de
ensaio, e apenas trocando o passo da régua, obtém-se a Tabela~\ref{tab:degradacao}.

\begin{table}[htbp]
\centering
\caption{Propagação do passo da régua para o menor degrau de \emph{slack} resolvível.
Valores em contagens são os medidos na versão Artix-7; a conversão para tempo usa o
passo de fase de cada dispositivo.}
\label{tab:degradacao}
\small
\begin{tabularx}{\textwidth}{@{}L{5.0cm} C{3.2cm} C{3.2cm} C{2.0cm}@{}}
\toprule
\textbf{Grandeza} & \textbf{Artix-7 (MMCM)} & \textbf{MAX 10 (PLL)} & \textbf{Razão} \\
\midrule
Passo de fase & \SI{19.8}{\pico\second} & \SI{96.2}{\pico\second} & $4{,}86\times$ \\
$\sigma_\text{resid}$, banco PID (\SI{0.726}{cnt}) & \SI{14.4}{\pico\second} & \SI{69.8}{\pico\second} & $4{,}86\times$ \\
$\sigma_\text{resid}$, banco \emph{bang-bang} (\SI{0.904}{cnt}) & \SI{17.9}{\pico\second} & \SI{86.9}{\pico\second} & $4{,}86\times$ \\
Degrau detectável $3\sigma$, PID & \SI{43}{\pico\second} & \SI{209}{\pico\second} & $4{,}86\times$ \\
Degrau detectável $3\sigma$, \emph{bang-bang} & \SI{54}{\pico\second} & \SI{261}{\pico\second} & $4{,}86\times$ \\
\bottomrule
\end{tabularx}
\end{table}

Sob o critério $\Delta t_\text{age} \geq k\,\sigma_\text{resid}$ que define
acionabilidade para SLM, o sensor migrado só é útil para decisões que tolerem
degraus de \emph{slack} da ordem de \SI{200}{\pico\second}. Se a campanha de estresse
viável não produzir degradação permanente dessa magnitude, o experimento não terá
poder de detecção --- e isso precisa ser verificado por cálculo, na Fase~0, antes de
qualquer investimento em implementação.

\subsection{Correspondência de recursos entre as plataformas}

A Tabela~\ref{tab:correspondencia} consolida o mapeamento de recursos. Ela serve de
referência rápida durante a Fase~0 e de índice para as frentes de implementação.

\begin{table}[htbp]
\centering
\caption{Correspondência de recursos entre a plataforma de referência e a plataforma alvo.}
\label{tab:correspondencia}
\small
\begin{tabularx}{\textwidth}{@{}L{3.0cm} L{3.4cm} X C{1.3cm}@{}}
\toprule
\textbf{Recurso} & \textbf{Artix-7 / Vivado} & \textbf{MAX 10 / Quartus Prime} & \textbf{Frente} \\
\midrule
Régua de fase &
MMCM, passo $T_\text{VCO}/56$ &
PLL, passo $T_\text{VCO}/8$; piso de \SI{96}{\pico\second}. \textbf{Perda de $4{,}86\times$.} & F1 \\
Protocolo de fase &
\texttt{PSEN}, \texttt{PSINCDEC}, \texttt{PSDONE} &
\texttt{phasestep}, \texttt{phaseupdown}, \texttt{phasecounterselect}, \texttt{phasedone}.
Correspondência praticamente um-para-um. & F1 \\
Recurso de atraso fino &
Cadeia \texttt{CARRY4}, \SI{\sim20}{\pico\second} por \emph{tap} &
Cadeia de \emph{carry} do LE de 4 entradas, dezenas a centena de ps por estágio. & F1 \\
Preservação lógica &
\texttt{DONT\_TOUCH}, \texttt{KEEP} &
Atributos \texttt{preserve}, \texttt{noprune}, \texttt{keep}, \texttt{dont\_merge}. & F2 \\
Posicionamento &
Restrições \texttt{LOC} e \texttt{BEL} &
\texttt{set\_location\_assignment LC\_X..Y..N..} no arquivo \texttt{.qsf}. & F2 \\
Inspeção de \emph{layout} &
Device View do Vivado &
Chip Planner. & F2 \\
Temperatura \emph{on-die} &
XADC, canal dedicado &
ADC SAR de 12~bits em modo diodo térmico (TSD), via IP \emph{Modular ADC Core}.
Exatidão inferior; exige calibração. & F3 \\
Tensão de núcleo &
XADC, monitor interno de $V_\text{CCINT}$ &
\textbf{Inexistente.} Requer derivação da trilha para pino ANAIN --- modificação de placa. & F3 \\
Transporte de dados &
Ponte USB-UART integrada (FT2232) &
Na DE10-Lite, o FTDI atende apenas JTAG. Alternativas: JTAG UART ou módulo
USB-serial externo no cabeçalho GPIO. & F4 \\
Supervisório embarcado &
MicroBlaze (opcional) &
Nios~II/e ou Nios~V/m, disponíveis sem licença na edição Lite. & F4 \\
Memória de configuração &
SRAM volátil, externa ao \emph{die} &
\textbf{Flash no mesmo \emph{die}} (CFM/UFM). Retenção é função de $T_j$ --- risco
específico sob estresse térmico prolongado. & F5 \\
\bottomrule
\end{tabularx}
\end{table}

\subsection{Riscos específicos do dispositivo alvo}

Três diferenças do MAX~10 não têm análogo na plataforma de referência e precisam de
tratamento explícito no plano.

\subsubsection{Memória de configuração no mesmo \emph{die}}

O Artix-7 carrega sua configuração em SRAM volátil a partir de memória externa. O
MAX~10 armazena a configuração em \emph{flash} integrada ao mesmo \emph{die} do tecido
lógico (blocos CFM e UFM). Consequência: durante uma campanha de burn-in, a memória de
configuração está sendo submetida ao mesmo estresse térmico do circuito sob teste, e as
especificações de retenção de dados da \emph{flash} são função da temperatura de junção.

Implicação prática: o envelope de estresse não pode ser definido apenas pelos limites
do tecido lógico. É necessário consultar a especificação de retenção do dispositivo,
adotar um setpoint conservador e prever verificação de integridade da configuração ao
final de cada campanha.

\subsubsection{Ausência de monitor interno de alimentação}

A perda do monitor de $V_\text{CCINT}$ é a mais consequente para a metodologia, porque
o eixo de tensão participa de duas funções distintas na versão de referência:
como covariável na decomposição de variância PVT, e como variável manipulada no
\emph{trim} de tensão em malha fechada.

A primeira função é recuperável com derivação da trilha de núcleo para um pino de
entrada analógica. A segunda exige interromper o regulador integrado da placa para
alimentar o núcleo pela fonte programável --- intervenção invasiva, irreversível e de
risco alto. O plano trata as duas separadamente e adia a segunda.

\subsubsection{Envelhecimento mais lento em nó tecnológico maior}

O MAX~10 é fabricado em \SI{55}{\nano\meter}. Os mecanismos dominantes de degradação
(NBTI e HCI) produzem, nesse nó, deslocamentos de $V_\text{th}$ menores e mais lentos
do que em \SI{28}{\nano\meter} planar, para o mesmo envelope de estresse. Combinado com
o piso de ruído cinco vezes mais alto, isso comprime a janela de detectabilidade pelos
dois lados simultaneamente.

\begin{warnbox}[title={Verificação de viabilidade obrigatória na semana 4}]
A trilha de estudo E1 deve responder, com número, à seguinte pergunta antes de qualquer
trabalho de implementação: \emph{quantas horas de estresse, no envelope térmico
admissível para este dispositivo, são necessárias para produzir um degrau de
\emph{slack} permanente de aproximadamente \SI{210}{\pico\second}?}

Se a resposta exceder a duração de campanha viável no laboratório, o projeto precisa
ser reformulado --- possivelmente como estudo de limite de detectabilidade, ou com
mudança de eixo de estresse. Esta decisão pertence à semana 4, não à semana 14.
\end{warnbox}

%=========================================================================
\section{Fase 0 --- Estudo}
%=========================================================================

\subsection{Justificativa e método}

A Fase~0 é obrigatória e bloqueante. Nenhum integrante abre o Quartus antes de concluir
a trilha de estudo correspondente à sua frente. A razão não é pedagógica em sentido
abstrato: em migrações desse tipo, o custo de descobrir tardiamente uma premissa
violada é medido em semanas de forno, e a maior parte dessas premissas está no
material de estudo.

\campo{Formato operacional}
\begin{itemize}
  \item Seminário semanal de 90 minutos, com um apresentador por trilha e arguição
        aberta ao restante da equipe.
  \item Cada trilha produz um artefato escrito de 2 a 4 páginas, versionado no
        repositório do projeto. O artefato não é resumo de leitura: é documento
        de projeto que será consumido pelas frentes de implementação.
  \item O orientador arguiu cada apresentador contra o critério de aceite declarado
        na trilha. Trilha não aprovada é reapresentada na semana seguinte.
  \item As seis trilhas correm em paralelo ao longo de quatro semanas. Um integrante
        é responsável primário por uma trilha e leitor secundário de outra.
\end{itemize}

\campo{Duração}
Quatro semanas, sobrepostas à preparação de bancada, com marco de saída na semana~4.

\subsection{Trilha E1 --- Física do envelhecimento e aceleração}

\campo{Conteúdo}
\begin{itemize}
  \item Mecanismos dominantes: NBTI (\emph{Negative Bias Temperature Instability})
        e HCI (\emph{Hot Carrier Injection}); origem física do deslocamento de
        tensão de limiar $V_\text{th}$.
  \item Dependência temporal em lei de potência; dependência térmica de Arrhenius;
        definição de fator de aceleração (AF) e de energia de ativação $E_a$;
        sensibilidade do AF à escolha de $E_a$.
  \item Distinção entre componente reversível (\emph{recovery}, variação PVT) e
        componente permanente do sinal --- e por que essa distinção é o problema
        central de medição do projeto.
  \item Por que o envelhecimento é observado \emph{indiretamente}, como perda gradual
        de margem temporal em lógica síncrona, em vez de medido diretamente.
  \item Conversão entre horas de estresse acelerado e horas equivalentes de operação
        em campo.
\end{itemize}

\campo{Leitura dirigida}
\begin{itemize}
  \item JEDEC JESD22-A108 --- \emph{Temperature, Bias and Operating Life}.
  \item JEDEC JESD91 --- métodos de extração de fator de aceleração.
  \item Artigos do grupo em SBCCI 2025 (validação do sensor) e SBCCI 2026
        (arquitetura de controle térmico e seu efeito sobre o AF entregue).
  \item Literatura de sensores in-situ de margem temporal: família Razor, canários
        de atraso, monitores baseados em oscilador em anel.
  \item Um artigo de revisão sobre SLM, para enquadrar o uso final do sensor
        (AVS, estimação de vida útil remanescente, manutenção preditiva).
\end{itemize}

\campo{Entregável}
Nota técnica derivando $\text{AF}(T,V)$ para os parâmetros da bancada e para o
dispositivo alvo, contendo obrigatoriamente a resposta numérica à pergunta de
viabilidade da Seção~2.5.3, com a faixa de $E_a$ considerada e a sensibilidade do
resultado a essa escolha.

\campo{Critério de aceite}
O apresentador explica, sem consulta, por que a redução da frequência de \emph{clock}
não corrige violação de \emph{hold}, e por que essa assimetria importa para o desenho
do caminho sintético sob teste.

\subsection{Trilha E2 --- Temporização, \emph{slack} e a régua de fase}

\campo{Conteúdo}
\begin{itemize}
  \item Fundamentos: tempo de \emph{setup}, tempo de \emph{hold}, atraso
        \emph{clock}-para-$Q$; caminho crítico e caminho curto; definição de
        \emph{slack} de \emph{setup}.
  \item Assimetria estrutural: \emph{setup} depende do período, \emph{hold} não.
        Correções distintas para cada violação.
  \item \emph{Skew} de \emph{clock} e seu efeito oposto sobre as duas restrições.
  \item Metaestabilidade e sincronizadores --- relevante porque o evento medido pelo
        sensor é, por construção, uma amostragem em região metaestável.
  \item Mecanismo central do projeto: como o deslocamento incremental de fase entre
        \emph{clock} de lançamento e \emph{clock} de captura converte um evento
        binário (falhou / não falhou) em uma medida temporal.
  \item Curva característica de probabilidade de falha em função da fase; a largura
        da transição como estimador de ruído; extração de precisão abaixo do passo
        por promediação.
  \item Por que uma régua travada em cristal é insensível a PVT e uma linha de atraso
        em tecido lógico não é.
\end{itemize}

\campo{Leitura dirigida}
\begin{itemize}
  \item \emph{MAX 10 Clocking and PLL User Guide}: capítulos de deslocamento de fase
        programável, reconfiguração dinâmica e faixa de operação do VCO.
  \item Documentação do MMCM da família 7 (UG472), para comparação lado a lado.
  \item Literatura de conversores tempo-digital (TDC): arquiteturas de linha de
        atraso, Vernier e batimento de frequência.
\end{itemize}

\campo{Entregável}
\begin{enumerate}
  \item Diagrama de tempos anotado do sensor atual, com identificação explícita da
        grandeza medida e do papel de cada borda.
  \item Tabela de correspondência de portas entre o protocolo de deslocamento de fase
        do MMCM e o do PLL do MAX~10, incluindo o \emph{handshake} de conclusão.
  \item Cálculo do passo de fase resultante para pelo menos três configurações de VCO,
        demonstrando o efeito da escolha.
\end{enumerate}

\campo{Critério de aceite}
O apresentador responde corretamente qual é o passo de fase obtido com VCO a
\SI{700}{\mega\hertz} versus \SI{1300}{\mega\hertz}, e explica por que a ferramenta
pode escolher o VCO automaticamente e como forçá-la a não fazê-lo.

\subsection{Trilha E3 --- O laço de auto-\emph{tuning}}

Esta é a trilha com maior densidade de informação não documentada, porque a fonte
primária é o próprio código.

\campo{Conteúdo}
\begin{itemize}
  \item Leitura dirigida do RTL atual, linha a linha, sob supervisão.
  \item O ciclo de medida: o que a máquina de estados executa entre duas amostras
        consecutivas de \emph{slack}.
  \item Inicialização: como a busca de borda é semeada na primeira medida, e o que
        acontece se a semente estiver fora da faixa útil.
  \item Rastreamento: como o laço acompanha o deslocamento da borda causado por
        variação térmica, variação de tensão e degradação permanente --- e como
        o laço, por construção, não distingue as três.
  \item Histerese e política de reaquisição em caso de perda de rastreamento.
  \item Comportamento nas extremidades da faixa de fase (\emph{wrap-around} ou
        saturação) e suas consequências para a série temporal.
  \item Parâmetros de sintonia expostos, faixas válidas e acoplamentos entre eles.
\end{itemize}

\campo{Entregável}
Diagrama de estados redesenhado \emph{a partir do código}, não a partir do artigo,
com anotação de todos os parâmetros de sintonia, suas faixas válidas e suas dependências.
Este documento passa a ser a especificação normativa para a reimplementação na frente F1.

\campo{Critério de aceite}
O apresentador identifica ao menos duas premissas do laço que dependem de o passo de
fase ser fino, e propõe reformulação de cada uma para operação com passo grosso.
Exemplos esperados: dimensionamento da histerese em número de passos, e critério
de convergência da busca inicial.

\subsection{Trilha E4 --- Instrumentação PVT}

\campo{Conteúdo}
\begin{itemize}
  \item O que o XADC entrega na versão de referência: canais, sequenciamento, taxa
        efetiva, ruído, e como esses dados alimentam a malha de compensação de tensão.
  \item Bloco ADC do MAX~10: conversor SAR de 12~bits, instanciação pelo IP
        \emph{Modular ADC Core} no Platform Designer, arquitetura do sequenciador,
        seleção de referência e faixa de entrada.
  \item Modo de diodo térmico (TSD): princípio, exatidão nominal, necessidade de
        calibração, e restrição de multiplexação com os canais analógicos externos.
  \item Ausência de monitor interno de trilha de alimentação e as consequências
        metodológicas discutidas na Seção~2.5.2.
  \item Propagação de incerteza: como o erro de temperatura entra no cálculo de AF.
\end{itemize}

\campo{Leitura dirigida}
\emph{MAX 10 Analog to Digital Converter User Guide}; folha de dados do dispositivo,
seções de características do ADC e do sensor térmico; documentação do XADC (UG480)
para comparação.

\campo{Entregável}
\begin{enumerate}
  \item Especificação da cadeia de medição de tensão de núcleo: ponto de acesso na
        placa, condicionamento, faixa, e orçamento de erro fechado.
  \item Procedimento de calibração do TSD contra referência externa, com número de
        pontos, faixa e método de ajuste.
  \item Estimativa de incerteza residual de temperatura após calibração, e sua
        propagação para o AF calculado na trilha E1.
\end{enumerate}

\campo{Critério de aceite}
O apresentador declara se a exatidão de temperatura alcançável é suficiente para o
cálculo de AF pretendido, com justificativa numérica, e não com juízo qualitativo.

\subsection{Trilha E5 --- Bancada de estresse}

\campo{Conteúdo}
\begin{itemize}
  \item Modelo FOPDT do forno; identificação por resposta ao degrau; sintonia por
        regras SIMC.
  \item Por que o controle \emph{bang-bang} subentrega estresse integrado a
        setpoint comandado idêntico --- o resultado central do artigo SBCCI 2026.
  \item Acoplamento térmico-elétrico do DUT: $T_\text{DUT} \rightarrow R_\text{DUT}
        \rightarrow I_\text{DUT} \rightarrow P_\text{DUT} \rightarrow T_\text{DUT}$,
        condições de ganho de laço e risco de fuga térmica.
  \item Por que a variável controlada relevante é a temperatura de junção do
        dispositivo, e não a temperatura do forno.
  \item Operação da fonte programável por VISA; limites de corrente; comportamento
        em transitório.
  \item Requisitos de JESD22-A108 quanto a monitoramento e auditabilidade das
        condições de estresse.
\end{itemize}

\campo{Entregável}
Plano de adaptação mecânica e térmica da placa alvo ao forno: fixação, roteamento de
cabos, posicionamento do termopar de referência, e caminho de saída dos sinais de
telemetria.

\campo{Critério de aceite}
O apresentador explica por que a sintonia do PID atual não é transferível para o novo
DUT, e quais grandezas físicas mudaram para justificar a reidentificação.

\subsection{Trilha E6 --- Aquisição e análise}

\campo{Conteúdo}
\begin{itemize}
  \item Leitura do software \emph{host} atual: protocolo serial, formato de pacote,
        política de \emph{timestamp}, sincronização entre fontes de dados
        heterogêneas, tratamento de reconexão.
  \item Requisitos de robustez para campanhas de vários dias: detecção de lacunas,
        integridade do arquivo em caso de queda, retomada.
  \item Pipeline de análise offline: decomposição de variância PVT do sinal de
        \emph{slack}, remoção de tendência, estimação de $\sigma_\text{resid}$,
        ajuste de modelos de degradação.
  \item Proveniência de dados: como um número de artigo é rastreado até o arquivo
        bruto que o gerou.
\end{itemize}

\campo{Entregável}
Ambiente de análise reproduzido localmente por cada integrante da frente, executando
sobre um CSV histórico do Artix-7 e reproduzindo os valores publicados dentro do
arredondamento.

\campo{Critério de aceite}
Reprodução numérica confirmada, com \emph{diff} apresentado contra os valores dos artigos.

\subsection{Marco de saída --- Reprodução de referência}

\begin{keybox}[title={Marco bloqueante da semana 4}]
Antes de qualquer trabalho sobre MAX~10, a equipe executa uma campanha curta
(2 a 4 horas) na Nexys~4~DDR existente, do início ao fim, com as próprias mãos:
síntese, gravação, montagem no forno, comando da fonte, aquisição, análise e relatório.
\end{keybox}

Este marco parece redundante e não é. Ele cumpre três funções que não podem ser
cumpridas de outra forma:

\begin{enumerate}
  \item \textbf{Produz o \emph{golden reference}.} Toda medida obtida em MAX~10 será
        comparada contra um conjunto de dados que a própria equipe gerou, com o
        instrumento em condição conhecida. Sem isso, discrepâncias futuras não têm
        atribuição possível.
  \item \textbf{Separa falhas de bancada de falhas de port.} Problemas de forno, de
        fonte, de cabo e de aquisição aparecem na semana 4, isoladamente. Sem esse
        marco, aparecem na semana 12, sobrepostos aos problemas de PLL, e o
        \emph{debug} conjunto custa múltiplas semanas.
  \item \textbf{Calibra a expectativa da equipe.} O contato físico com o sistema
        funcionando estabelece o padrão de qualidade contra o qual o port será julgado.
\end{enumerate}

\subsection{Matriz de cobertura}

A Tabela~\ref{tab:cobertura} relaciona trilhas de estudo e frentes de implementação.
Ela é usada para verificar que nenhuma frente inicia sem sua base conceitual e para
alocar leitores secundários.

\begin{table}[htbp]
\centering
\caption{Matriz trilha de estudo $\times$ frente de implementação.
$\bullet$ indica pré-requisito bloqueante; $\circ$ indica leitura recomendada.}
\label{tab:cobertura}
\small
\begin{tabular}{@{}lccccccl@{}}
\toprule
\textbf{Trilha} & \textbf{F1} & \textbf{F2} & \textbf{F3} & \textbf{F4} & \textbf{F5} & \textbf{F6} \\
\midrule
E1 --- Física do envelhecimento & $\circ$ & & & & $\bullet$ & $\bullet$ \\
E2 --- Temporização e régua & $\bullet$ & $\bullet$ & & & & $\circ$ \\
E3 --- Laço de auto-\emph{tuning} & $\bullet$ & $\circ$ & & $\circ$ & & $\circ$ \\
E4 --- Instrumentação PVT & & & $\bullet$ & $\bullet$ & $\circ$ & $\circ$ \\
E5 --- Bancada de estresse & & & $\circ$ & & $\bullet$ & $\circ$ \\
E6 --- Aquisição e análise & & & & $\bullet$ & & $\bullet$ \\
\bottomrule
\end{tabular}
\end{table}

\clearpage
%=========================================================================
\section{Frentes de implementação}
%=========================================================================

\subsection{Princípio de decomposição e alocação}

O trabalho é decomposto em seis frentes, definidas por acoplamento técnico e não por
conveniência de agenda. A regra de alocação é:

\begin{itemize}
  \item Dimensionamento nominal: 4 a 6 integrantes, um responsável primário por frente,
        com frentes de menor carga acumuladas por um mesmo integrante.
  \item Com equipe reduzida, F2 colapsa em F1 e F6 colapsa em F4.
  \item \textbf{F1 nunca compartilha responsável com outra frente.} É a frente crítica,
        determina a viabilidade do projeto e não tolera divisão de atenção.
  \item Toda frente tem um segundo integrante designado como revisor, que não executa
        mas responde pela conferência dos entregáveis.
\end{itemize}

\begin{table}[htbp]
\centering
\caption{Visão geral das frentes.}
\label{tab:frentes}
\small
\begin{tabularx}{\textwidth}{@{}L{1.0cm} L{4.3cm} X C{1.9cm}@{}}
\toprule
\textbf{ID} & \textbf{Frente} & \textbf{Questão que responde} & \textbf{Criticidade} \\
\midrule
F1 & Régua temporal e RTL do sensor & Quão fina e quão estável é a referência de tempo disponível? & Crítica \\
F2 & Implementação física e fluxo Quartus & O circuito descrito sobrevive à ferramenta e se repete entre compilações? & Alta \\
F3 & Instrumentação PVT e hardware de placa & As covariáveis de temperatura e tensão estão disponíveis e calibradas? & Alta \\
F4 & Telemetria e cadeia de aquisição & Os dados chegam íntegros e sincronizados após dias de campanha? & Média \\
F5 & Bancada térmica e elétrica & O estresse é entregue de forma controlada e sem destruir o dispositivo? & Média \\
F6 & Caracterização, dados e produção científica & Quanto vale a medida obtida, com que incerteza? & Alta \\
\bottomrule
\end{tabularx}
\end{table}

%-------------------------------------------------------------------------
\subsection{F1 --- Régua temporal e RTL do sensor}

\frenteheader{F1 --- Régua temporal e RTL do sensor}
{Estabelecer no MAX 10 uma referência de tempo rastreável e tão fina quanto possível, e reimplementar sobre ela o laço de auto-tuning.}
{Frente crítica. Determina a viabilidade do projeto e o escopo de tudo o que vem depois.}

\campo{Pré-requisitos}
Trilhas E2 e E3 aprovadas. Acesso ao RTL da versão de referência e ao diagrama de
estados produzido em E3.

\campo{Tarefas}
\begin{enumerate}
  \item \textbf{Instanciação e caracterização do PLL.}
        Instanciar PLL com reconfiguração dinâmica de fase (ALTPLL com reconfiguração
        habilitada, ou Altera PLL acompanhado do Altera PLL Reconfig). Forçar o VCO ao
        teto da faixa admissível, verificando o valor efetivamente escolhido pela
        ferramenta e não o valor pedido. Medir o passo de fase real, sem confiar na
        especificação: montar experimento de contagem sobre caminho de atraso conhecido
        e levantar o passo médio e sua dispersão.

  \item \textbf{Porte da máquina de estados.}
        Traduzir o protocolo de \emph{handshake} de deslocamento de fase, com atenção
        particular ao sinal de conclusão --- a semântica de \texttt{phasedone} difere
        da de \texttt{PSDONE} em detalhes de temporização que precisam ser verificados
        em simulação antes de irem para \emph{hardware}. Reformular os parâmetros
        identificados em E3 como dependentes de passo fino.

  \item \textbf{Redesenho do caminho sintético sob teste.}
        Reescrever o caminho crítico para o tecido lógico de 4 entradas do MAX~10, com
        atraso alvo dimensionado para a nova granularidade: o caminho precisa ser longo
        o suficiente para que a borda de falha caia em região de fase bem resolvida, e
        curto o suficiente para caber na faixa de deslocamento disponível.

  \item \textbf{Investigação de Vernier por batimento entre dois PLLs.}
        Esta é a tarefa de maior valor e maior risco da frente. O dispositivo dispõe de
        múltiplos PLLs com possibilidade de cascateamento por GCLK. Dois \emph{clocks}
        com diferença de frequência pequena fornecem resolução temporal efetiva
        $1/f_1 - 1/f_2$, arbitrariamente fina, ao custo de tempo de medição por amostra.
        É a única rota que recupera resolução \emph{mantendo a régua travada em cristal}.
        \begin{itemize}
          \item Modelar a resolução efetiva e o tempo de medida em função de $\Delta f$.
          \item Verificar se as restrições de contadores do dispositivo permitem os
                pares de frequência necessários.
          \item Avaliar o impacto do \emph{jitter} relativo entre os dois PLLs sobre a
                resolução efetivamente utilizável --- este é o fator que pode invalidar
                a abordagem, e precisa ser medido, não estimado.
          \item Prototipar e caracterizar.
        \end{itemize}

  \item \textbf{Rota alternativa: interpolação em tecido lógico.}
        Executada apenas se a tarefa 4 for declarada inviável no Gate~A. Subdivisão do
        passo de \SI{96}{\pico\second} por cadeia de atraso em células lógicas, com
        calibração por densidade de código em cada ponto de temperatura da campanha.
        Se adotada, a limitação de rastreabilidade deve ser declarada explicitamente
        em todo relatório e publicação derivados --- a régua passa a variar com PVT
        e a envelhecer junto com o mensurando.

  \item \textbf{Banco de testes.}
        Simulação do laço completo com modelo de deslocamento de borda, incluindo os
        casos de extremidade de faixa e de perda de rastreamento.
\end{enumerate}

\campo{Entregáveis}
\begin{itemize}
  \item RTL parametrizável, com passo de fase e faixa como parâmetros de topo.
  \item Banco de testes de simulação cobrindo inicialização, rastreamento, extremidade
        de faixa e reaquisição.
  \item Relatório de caracterização da régua: passo médio medido, dispersão,
        não-linearidade diferencial, deriva com temperatura.
  \item Parecer técnico sobre a viabilidade do Vernier, com dados de suporte.
\end{itemize}

\campo{Critério de aceite --- Gate A}
Passo de fase medido, documentado e estável, e laço de auto-\emph{tuning} rastreando
uma borda deslocada artificialmente por variação controlada de temperatura ambiente,
em \emph{hardware}.

\campo{Riscos}
\begin{itemize}
  \item \emph{Jitter} relativo entre PLLs anula o ganho teórico do Vernier.
  \item Restrições de contadores impedem os pares de frequência desejados.
  \item Tempo de medida por amostra do Vernier torna a taxa de amostragem incompatível
        com a dinâmica térmica que se deseja observar.
\end{itemize}

\campo{Perfil do integrante}
Maior maturidade em RTL e, principalmente, em leitura crítica de documentação de
dispositivo. Alocar o integrante mais forte da equipe. Tarefa adequada a aluno de
pós-graduação.

%-------------------------------------------------------------------------
\subsection{F2 --- Implementação física e fluxo Quartus}

\frenteheader{F2 --- Implementação física e fluxo Quartus}
{Garantir que o circuito descrito em RTL sobreviva às otimizações da ferramenta e que compilações sucessivas produzam o mesmo circuito físico.}
{Frente habilitadora. Sem reprodutibilidade de implementação, nenhuma medida é comparável.}

\campo{Pré-requisitos}
F1 com RTL preliminar disponível. Trilha E2 aprovada.

\campo{Tarefas}
\begin{enumerate}
  \item \textbf{Tradução do vocabulário de restrições.}
        Substituir as diretivas de preservação e posicionamento do fluxo Xilinx pelos
        equivalentes Quartus: atributos \texttt{preserve}, \texttt{noprune},
        \texttt{keep} e \texttt{dont\_merge} no código; atribuições
        \texttt{set\_location\_assignment LC\_X..Y..N..} no arquivo \texttt{.qsf} para
        fixação de células.

  \item \textbf{Desabilitação seletiva de otimizações.}
        Identificar e desligar explicitamente as transformações do \emph{fitter} que
        reescrevem o caminho sintético: \emph{retiming} de registradores, duplicação
        de registradores, compartilhamento de recursos, remoção de lógica redundante.
        \textbf{Cada atribuição deve ser comentada no \texttt{.qsf} com sua justificativa},
        porque este arquivo será lido por terceiros ao avaliar o trabalho.

  \item \textbf{Instanciação direta de célula do dispositivo, se necessário.}
        Caso a inferência não produza a estrutura desejada, instanciar diretamente a
        célula lógica WYSIWYG da família, com portas de \emph{carry} de entrada e saída.
        \emph{Confirmar o nome exato da primitiva na biblioteca do dispositivo dentro da
        versão instalada do Quartus} --- não em documentação de terceiros nem em fórum.

  \item \textbf{Estabelecimento de reprodutibilidade.}
        Definir e verificar o procedimento que garante posicionamento idêntico entre
        compilações: controle de \emph{seed}, verificação de \emph{netlist} pós-síntese,
        conferência visual no Chip Planner, e comparação automatizada de relatórios
        de posicionamento.

  \item \textbf{Restrições de temporização.}
        Escrever o arquivo SDC com as exceções corretas para o caminho sob medição ---
        que é, por construção, um caminho que se pretende violar. Documentar por que
        cada exceção existe.

  \item \textbf{Análise de recursos e de \emph{floorplan}.}
        Definir a região do dispositivo ocupada pelo sensor e sua vizinhança lógica,
        de modo a manter constante o ambiente térmico e de acoplamento entre versões.
\end{enumerate}

\campo{Entregáveis}
\begin{itemize}
  \item Projeto Quartus completo com \texttt{.qsf} integralmente comentado.
  \item Roteiro escrito de verificação de reprodutibilidade.
  \item Relatório de posicionamento com capturas do Chip Planner e comparação
        entre compilações.
  \item Arquivo SDC documentado.
\end{itemize}

\campo{Critério de aceite}
Dez compilações consecutivas, a partir do mesmo código-fonte, produzem posicionamento
idêntico do bloco sensor, verificado por comparação automatizada.

\campo{Riscos}
\begin{itemize}
  \item Otimização não documentada da ferramenta alterando o caminho entre versões
        do Quartus --- mitigado fixando a versão da ferramenta para todo o projeto.
  \item Restrições de posicionamento conflitantes tornando o \emph{fit} inviável.
\end{itemize}

\campo{Perfil do integrante}
Paciência com ferramenta e disposição para leitura de relatórios de compilação.
Adequado a aluno de graduação com experiência prévia em FPGA.

%-------------------------------------------------------------------------
\subsection{F3 --- Instrumentação PVT e hardware de placa}

\frenteheader{F3 --- Instrumentação PVT e hardware de placa}
{Recuperar os dois eixos de covariável --- temperatura de junção e tensão de núcleo --- que o XADC fornecia nativamente na plataforma de referência.}
{Frente com intervenção física irreversível. Exige revisão de projeto antes da execução.}

\campo{Pré-requisitos}
Trilha E4 aprovada, incluindo o orçamento de erro.

\campo{Tarefas}
\begin{enumerate}
  \item \textbf{Instanciação e caracterização do ADC.}
        Configurar o IP \emph{Modular ADC Core} no Platform Designer. Caracterizar taxa
        efetiva de amostragem, ruído de conversão, comportamento do sequenciador e
        latência de leitura.

  \item \textbf{Modo de diodo térmico e calibração.}
        Colocar o modo TSD em operação. Calibrar contra termopar de referência em, no
        mínimo, cinco pontos distribuídos pela faixa de estresse prevista, com o
        dispositivo em regime permanente em cada ponto. Levantar a curva de erro
        residual pós-calibração e sua dependência com a temperatura.

  \item \textbf{Esquema de intercalação TSD / canais externos.}
        O modo de sensoriamento térmico e a leitura de canais analógicos externos não
        são simultâneos. Definir e documentar o esquema de alternância temporal adotado,
        e quantificar o erro de sincronismo que ele introduz entre as séries de
        temperatura e de tensão.

  \item \textbf{Acesso à tensão de núcleo.}
        Levantar o ponto de acesso à trilha de núcleo na placa. Projetar a derivação
        para pino de entrada analógica, incluindo condicionamento e proteção.
        \begin{itemize}
          \item Revisão de projeto obrigatória com o orientador antes da execução.
          \item Documentação fotográfica antes, durante e depois.
          \item Procedimento de recuperação escrito antes da intervenção.
        \end{itemize}

  \item \textbf{Integração com a cadeia de telemetria.}
        Entregar as leituras de temperatura e tensão à frente F4 no mesmo pacote e com
        o mesmo \emph{timestamp} da medida de \emph{slack}.
\end{enumerate}

\begin{warnbox}[title={Decisão de escopo: eixo de estresse elétrico}]
Alimentar o núcleo do dispositivo pela fonte programável exige interromper o regulador
integrado da placa. É a intervenção de maior risco do projeto inteiro e a de menor
probabilidade de sucesso na primeira tentativa.

\textbf{Recomendação:} a fase~1 do projeto opera \emph{sem} estresse de tensão, com
monitoramento passivo de $V_\text{CORE}$ e temperatura como único eixo de estresse. O
eixo de tensão entra apenas na fase~2, após o Gate~A, e somente se houver segunda placa
disponível como sacrificial.

Consequência metodológica a ser declarada: a decomposição de variância PVT da fase~1
terá o eixo de tensão como covariável observada, não como variável manipulada. A
comparação com os resultados do Artix-7 deve ser restrita aos eixos comuns.
\end{warnbox}

\campo{Entregáveis}
\begin{itemize}
  \item Curva de calibração do TSD com incerteza declarada.
  \item Documentação esquemática e fotográfica da modificação de placa.
  \item Procedimento de recuperação e registro de intervenções.
  \item Especificação do formato de pacote de telemetria PVT.
\end{itemize}

\campo{Critério de aceite}
Leitura calibrada e simultânea de temperatura de junção e tensão de núcleo, transmitida
ao \emph{host} junto com a medida de \emph{slack}, com erro de sincronismo quantificado.

\campo{Perfil do integrante}
Habilidade de bancada e cuidado com \emph{hardware}. Experiência prévia em soldagem
fina é requisito, não diferencial.

%-------------------------------------------------------------------------
\subsection{F4 --- Telemetria e cadeia de aquisição}

\frenteheader{F4 --- Telemetria e cadeia de aquisição}
{Levar dados sincronizados e íntegros do dispositivo sob teste ao host, continuamente, por campanhas de vários dias.}
{Frente de infraestrutura. Falha aqui invalida semanas de forno.}

\campo{Pré-requisitos}
Trilhas E4 e E6 aprovadas.

\campo{Tarefas}
\begin{enumerate}
  \item \textbf{Decisão de transporte.}
        A placa alvo não expõe ponte USB-UART de uso geral: o conversor FTDI integrado
        atende apenas à programação JTAG. Duas alternativas:
        \begin{itemize}
          \item \emph{Módulo USB-serial externo} no cabeçalho GPIO, replicando o
                protocolo atual sem alteração no \emph{host}. Vantagem: compatibilidade
                imediata. Desvantagem: mais um componente e mais dois cabos atravessando
                a parede do forno.
          \item \emph{JTAG UART} da Intel, com terminal via System Console. Vantagem:
                telemetria pelo mesmo cabo da programação, sem \emph{hardware} adicional
                dentro do forno. Desvantagem: exige adaptação da camada de transporte
                no \emph{host}.
        \end{itemize}
        \textbf{Recomendação:} JTAG UART na fase~1, pela redução de \emph{hardware}
        exposto a temperatura; migração para UART externo apenas se a robustez em
        campanha longa se mostrar insuficiente.

  \item \textbf{Porte do protocolo de pacote.}
        Preservar o formato de saída CSV compatível com o \emph{pipeline} de análise
        existente. Esta compatibilidade é o que torna a comparação Artix-7 $\times$
        MAX~10 direta, e não deve ser sacrificada por conveniência de implementação.

  \item \textbf{Política de \emph{timestamp}.}
        Definir a fonte de tempo autoritativa e o mecanismo de correlação entre as
        séries de \emph{slack}, temperatura, tensão e telemetria da fonte. Quantificar
        o erro de correlação.

  \item \textbf{Adaptação do software \emph{host}.}
        Reconexão automática, detecção e marcação de lacunas, gravação incremental
        resistente a queda, gráficos ao vivo.

  \item \textbf{Ensaio de resistência.}
        Setenta e duas horas contínuas com injeção deliberada de falhas: desconexão
        de USB, suspensão do \emph{host}, saturação de disco.
\end{enumerate}

\campo{Entregáveis}
\begin{itemize}
  \item Camada de telemetria no DUT e no \emph{host}.
  \item Especificação escrita do protocolo e do formato CSV.
  \item Relatório do ensaio de \SI{72}{\hour} com estatística de perda de amostras e
        de lacunas detectadas.
\end{itemize}

\campo{Critério de aceite}
Setenta e duas horas contínuas com perda de amostras abaixo do limiar acordado e sem
perda de sincronismo de \emph{timestamp} em nenhuma das falhas injetadas.

\campo{Perfil do integrante}
Conforto com Python e com linha de comando; disposição para escrever código defensivo.

%-------------------------------------------------------------------------
\subsection{F5 --- Bancada térmica e elétrica}

\frenteheader{F5 --- Bancada térmica e elétrica}
{Adaptar a bancada de burn-in existente ao novo dispositivo sob teste, respeitando os limites específicos do MAX 10.}
{Frente com risco de perda de hardware. O envelope de estresse é decisão de projeto, não parâmetro livre.}

\campo{Pré-requisitos}
Trilhas E1 e E5 aprovadas. Marco de reprodução de referência concluído.

\campo{Tarefas}
\begin{enumerate}
  \item \textbf{Adaptação mecânica.}
        Fixação da placa no forno, roteamento de cabos de alimentação e telemetria,
        posicionamento do termopar de referência em contato controlado com o encapsulamento.

  \item \textbf{Reidentificação da planta.}
        A massa térmica, a área de troca e a impedância térmica mudaram com a troca de
        placa. Os ganhos do PID atual não são transferíveis. Executar nova identificação
        FOPDT por resposta ao degrau e nova sintonia por regras SIMC, documentando os
        parâmetros identificados.

  \item \textbf{Definição do envelope de estresse.}
        Fixar setpoint de temperatura e duração de campanha respeitando simultaneamente:
        \begin{itemize}
          \item o limite de temperatura de junção do dispositivo na versão comercial
                (\SI{85}{\celsius});
          \item a especificação de retenção da memória \emph{flash} de configuração
                integrada ao \emph{die}, que é função de $T_j$ e do tempo acumulado ---
                restrição \emph{sem análogo} na plataforma de referência;
          \item o resultado de viabilidade da trilha E1, que estabelece o piso de
                estresse necessário para produzir sinal detectável.
        \end{itemize}
        Se o piso da trilha E1 exceder o teto imposto pelas duas primeiras restrições,
        essa contradição é o resultado principal da fase e deve ser reportada como tal.

  \item \textbf{Verificação de integridade pós-campanha.}
        Procedimento de conferência da memória de configuração ao término de cada
        campanha: releitura, comparação com imagem de referência, registro do resultado.

  \item \textbf{Protocolo operacional de campanha.}
        Rampa de aquecimento, patamar, leituras periódicas com forno desligado para
        separação de componente reversível, resfriamento controlado, e registro
        auditável de todas as condições conforme JESD22-A108.
\end{enumerate}

\campo{Entregáveis}
\begin{itemize}
  \item Modelo FOPDT identificado e controlador PID sintonizado, com dados de validação.
  \item Documento de envelope de estresse com justificativa de cada limite.
  \item Procedimento operacional padrão de campanha.
  \item Registro de verificação de integridade de configuração.
\end{itemize}

\campo{Critério de aceite}
Rastreamento de setpoint dentro de orçamento de erro médio equivalente ao do banco PID
da plataforma de referência, demonstrado em campanha de validação de pelo menos
\SI{12}{\hour}.

\campo{Perfil do integrante}
Interesse em controle e instrumentação. Adequado a aluno com disciplina de sistemas de
controle cursada.

%-------------------------------------------------------------------------
\subsection{F6 --- Caracterização, dados e produção científica}

\frenteheader{F6 --- Caracterização, dados e produção científica}
{Transformar a implementação em resultado defensável, com incerteza declarada e método auditável.}
{Frente que define o valor final do projeto. Inicia cedo, não no fim.}

\campo{Pré-requisitos}
Trilha E6 aprovada e reprodução de referência concluída.

\campo{Tarefas}
\begin{enumerate}
  \item \textbf{Caracterização da régua.}
        Passo médio, dispersão, não-linearidade diferencial, deriva com temperatura.
        Esta caracterização é a base de toda conversão de contagens para picossegundos
        e precisa ser refeita se qualquer parâmetro de PLL mudar.

  \item \textbf{Decomposição de variância PVT.}
        Aplicar ao sinal de \emph{slack} do MAX~10 a mesma metodologia dos artigos
        anteriores, com atenção à mudança de status do eixo de tensão (covariável
        observada, não manipulada, na fase~1).

  \item \textbf{Estimação de $\sigma_\text{resid}$ e do degrau resolvível.}
        Cálculo em picossegundos, com barra de incerteza e método de estimação
        documentado. Declaração explícita do menor degrau de \emph{slack} resolvível
        sob critério $k\sigma$, para $k = 3$ e ao menos um valor alternativo.

  \item \textbf{Comparação sistemática entre plataformas.}
        Tabela de comparação em eixos comparáveis, com incertezas, e discussão
        explícita do que \emph{não} é comparável e por quê.

  \item \textbf{Gestão de dados e proveniência.}
        Versionamento dos conjuntos de dados brutos; \emph{notebook} de análise
        versionado; rastreabilidade de cada número reportado até o arquivo bruto
        que o originou.

  \item \textbf{Produção escrita.}
        Relatório técnico interno; material para submissão a SBCCI ou LATS sobre
        portabilidade entre fabricantes de sensores de \emph{slack}, incluindo
        o resultado desfavorável, se for esse o desfecho.
\end{enumerate}

\campo{Entregáveis}
\begin{itemize}
  \item \emph{Notebook} de análise versionado e reprodutível.
  \item Relatório de caracterização da régua.
  \item Tabela comparativa entre plataformas com incertezas.
  \item Rascunho de artigo.
\end{itemize}

\campo{Critério de aceite}
$\sigma_\text{resid}$ do MAX~10 reportado em picossegundos, com barra de incerteza e
método de estimação auditável por terceiro que não participou da medição.

\campo{Perfil do integrante}
Maturidade em análise de dados e em escrita técnica. Adequado a aluno de pós-graduação,
possivelmente o mesmo responsável por F1 em equipe reduzida.

\clearpage
%=========================================================================
\section{Sequenciamento, portões e cronograma}
%=========================================================================

\subsection{Cronograma de referência}

O cronograma abaixo assume um semestre letivo de 16 semanas e dedicação parcial dos
integrantes. É um cronograma de dependências, não de esforço: as datas dos portões são
firmes; as atividades entre portões absorvem variação.

\begin{table}[htbp]
\centering
\caption{Cronograma de referência com portões de decisão.}
\label{tab:cronograma}
\small
\begin{tabularx}{\textwidth}{@{}C{1.6cm} L{5.2cm} X@{}}
\toprule
\textbf{Semanas} & \textbf{Atividade principal} & \textbf{Marco / portão} \\
\midrule
1--4 & Fase 0: seis trilhas de estudo em paralelo; preparação de bancada em segundo plano &
Todas as trilhas aprovadas em arguição \\
4 & Reprodução de referência na Nexys 4 DDR &
\emph{Golden reference} depositado no repositório \\
5--8 & F1 e F2 em execução paralela; F3 e F5 em preparação; F4 definindo transporte &
\textbf{Gate A} (semana 8): régua caracterizada e laço rastreando em \emph{hardware} \\
9--12 & F3 e F4 em integração; primeira campanha curta de validação &
\textbf{Gate B} (semana 12): campanha de \SI{24}{\hour} com cadeia completa de dados \\
13--16 & F5 executa campanha longa; F6 conduz análise e redação &
Relatório técnico e rascunho de artigo \\
\bottomrule
\end{tabularx}
\end{table}

\subsection{Gate A --- Viabilidade da régua (semana 8)}

\begin{keybox}[title={Gate A é o portão real do projeto}]
Gate~A decide se o projeto continua sendo uma migração ou passa a ser outro projeto.
A decisão é do orientador, não da equipe, e é tomada na semana~8 com os dados
disponíveis na semana~8 --- não adiada para quando houver mais dados.
\end{keybox}

\campo{Evidência exigida}
\begin{enumerate}
  \item Passo de fase do PLL medido em \emph{hardware}, com dispersão.
  \item Laço de auto-\emph{tuning} rastreando borda deslocada por variação térmica.
  \item Parecer sobre o Vernier de dois PLLs, com dados de \emph{jitter} relativo.
\end{enumerate}

\campo{Desfechos possíveis}
\begin{itemize}
  \item \textbf{Vernier caracterizado e viável.} O projeto segue como migração com
        recuperação de resolução. Escopo original mantido. F1 conclui a implementação
        do Vernier nas semanas 9--12.
  \item \textbf{Vernier inviável, passo nativo funcionando.} O projeto é reposicionado
        como \emph{estudo de limite de detectabilidade de sensores de slack em nó
        tecnológico maior}. As metas de F6 mudam: o produto passa a ser a quantificação
        rigorosa da penalidade, e não a demonstração de equivalência. Continua sendo
        resultado publicável.
  \item \textbf{Laço não rastreia de forma estável.} Situação mais grave. Investigar se
        a causa é de F1 (reformulação de parâmetros) ou de F2 (fitter reescrevendo o
        caminho). Prorrogação máxima de duas semanas, com Gate~A reavaliado na semana 10
        e compressão proporcional da fase de campanha.
\end{itemize}

\subsection{Gate B --- Integridade da cadeia (semana 12)}

Verifica que a cadeia completa --- sensor, PVT, telemetria, \emph{host}, análise ---
opera em conjunto antes de comprometer semanas de forno.

\campo{Evidência exigida}
Campanha de \SI{24}{\hour} produzindo arquivo único, sem lacunas não marcadas, com
séries de \emph{slack}, temperatura e tensão correlacionadas, e análise preliminar
executada sobre o arquivo pela frente F6.

\campo{Critério de bloqueio}
Se a análise preliminar não consegue produzir uma estimativa de $\sigma_\text{resid}$
com o arquivo obtido, a campanha longa não é autorizada. O problema está na cadeia,
e uma campanha de duas semanas apenas produzirá o mesmo problema em maior volume.

\subsection{Ritual de acompanhamento}

\begin{itemize}
  \item \textbf{Reunião semanal de 60 minutos} com toda a equipe. Cada frente reporta:
        o que foi feito, o que está bloqueado, o que vai ser feito. Ata escrita,
        versionada, com decisões destacadas.
  \item \textbf{Revisão cruzada obrigatória.} Todo entregável passa pelo revisor
        designado da frente antes de ir ao orientador.
  \item \textbf{Repositório único} para código, dados, documentos e atas.
        Convenção de \emph{commits} adotada e verificada.
  \item \textbf{Caderno de bancada}, físico ou digital, para toda intervenção em
        \emph{hardware} e toda campanha. Data, condições, quem executou, o que
        deu errado.
\end{itemize}

%=========================================================================
\section{Gestão de risco}
%=========================================================================

\begin{table}[htbp]
\centering
\caption{Registro de riscos do projeto, ordenado por impacto.}
\label{tab:riscos}
\small
\begin{tabularx}{\textwidth}{@{}L{4.0cm} C{1.3cm} X@{}}
\toprule
\textbf{Risco} & \textbf{Frente} & \textbf{Mitigação} \\
\midrule
Sinal de envelhecimento em \SI{55}{\nano\meter} permanece abaixo do piso de ruído em qualquer campanha viável &
F5, F6 &
Trilha E1 responde por cálculo na semana 4, antes de qualquer investimento em implementação. Se confirmado, o projeto é reformulado imediatamente como estudo de limite de detectabilidade. \\
Vernier de dois PLLs inviável por \emph{jitter} relativo &
F1 &
Caracterização obrigatória até o Gate~A. Rota alternativa (interpolação em tecido) já especificada, com limitação de rastreabilidade declarada. Reposicionamento de escopo previsto. \\
Degradação da \emph{flash} de configuração por estresse térmico prolongado &
F5 &
Envelope de estresse limitado pela especificação de retenção; verificação de integridade da configuração ao fim de cada campanha; placa sacrificial designada. \\
Modificação de placa para acesso a $V_\text{CORE}$ danifica o \emph{hardware} &
F3 &
Revisão de projeto antes da execução; procedimento de recuperação escrito previamente; adiamento do eixo de tensão para a fase 2; segunda placa. \\
\emph{Fitter} reescrevendo o caminho sintético entre compilações &
F2 &
Critério de reprodutibilidade de dez compilações é bloqueante para o Gate~A. Versão do Quartus fixada para todo o projeto. \\
Perda de dados em campanha longa &
F4 &
Ensaio de resistência de \SI{72}{\hour} com injeção de falhas antes do Gate~B; gravação incremental; detecção e marcação de lacunas. \\
Erro de calibração do TSD propagando para o AF &
F3, F6 &
Orçamento de erro fechado na trilha E4; calibração em cinco pontos; propagação de incerteza explicitada em toda conclusão que dependa do AF. \\
Dispersão da equipe entre frentes, com F1 subatendida &
--- &
Regra de alocação exclusiva para F1; revisor designado por frente; reunião semanal com relato de bloqueios. \\
Perda de comparabilidade com a plataforma de referência por divergência de formato &
F4 &
Compatibilidade de formato CSV tratada como requisito, não como preferência. \\
\bottomrule
\end{tabularx}
\end{table}

%=========================================================================
\section{Infraestrutura, ferramentas e convenções}
%=========================================================================

\subsection{Ambiente de ferramenta}

\begin{itemize}
  \item \textbf{Versão do Quartus Prime fixada} no início do projeto e registrada no
        repositório. Otimizações de \emph{fitter} mudam entre versões, e uma
        atualização no meio da campanha invalida a comparabilidade entre medidas.
  \item Projeto versionado por arquivos-fonte (\texttt{.qsf}, \texttt{.sdc}, HDL, IP
        em formato textual), nunca por diretório de compilação.
  \item Ambiente Python do \emph{host} e da análise fixado por arquivo de dependências.
\end{itemize}

\subsection{Repositório e convenções}

\begin{itemize}
  \item Repositório único, com separação clara entre \texttt{hdl/}, \texttt{host/},
        \texttt{analysis/}, \texttt{data/}, \texttt{docs/} e \texttt{minutes/}.
  \item Convenção de \emph{commits} adotada e verificada em revisão.
  \item Dados brutos imutáveis: nenhum arquivo de campanha é editado após a gravação.
        Correções são feitas em camada de processamento, com registro.
\end{itemize}

\subsection{Proveniência de dados}

Todo número que apareça em relatório ou artigo precisa ser rastreável, por caminho
documentado, até o arquivo bruto que o originou e até a versão do código de análise
que o produziu. Esta exigência não é burocrática: em campanhas de vários dias com
múltiplas fontes de telemetria, a incapacidade de rastrear um número é indistinguível
de erro.

\subsection{Segurança de bancada}

\begin{itemize}
  \item Nenhuma campanha noturna sem supervisão remota de temperatura e sem limite
        de corrente configurado na fonte.
  \item Procedimento de desligamento de emergência afixado na bancada.
  \item Intervenções em \emph{hardware} sempre com segunda pessoa presente.
\end{itemize}

%=========================================================================
\section{Notas de orientação}
%=========================================================================

Três pontos merecem ser comunicados à equipe verbalmente, na primeira reunião, antes
da distribuição deste documento. Correspondem aos modos de falha mais frequentes em
projetos desta natureza conduzidos por equipes em formação.

\subsection{A régua não é o sensor}

A tentação natural do aluno é reimplementar a lógica, ver números saírem pela serial e
considerar o trabalho concluído. O objeto científico deste projeto é a \emph{incerteza
da medida}, e ela nasce inteiramente da referência de tempo. Um sensor que compila e
produz séries plausíveis, sem caracterização de régua, não constitui resultado ---
produz números cuja unidade é desconhecida.

Consequência prática para a orientação: cobrar a caracterização da régua \emph{antes}
de cobrar o funcionamento do laço, mesmo que a ordem pareça invertida.

\subsection{Tecido lógico não é referência}

Se a resolução for recuperada por linha de atraso em células lógicas, a régua passa a
variar com PVT e a envelhecer junto com o mensurando. Isso não invalida a abordagem,
mas transforma o problema: exige calibração por ponto de temperatura e introduz um
termo de erro que precisa ser propagado até a conclusão final.

O erro a evitar não é adotar essa rota --- é adotá-la sem declarar a limitação, ou
sem perceber que ela existe.

\subsection{Resultado desfavorável bem medido é resultado}

A afirmação ``a arquitetura não porta com fidelidade útil para este nó tecnológico,
pelas razões X e Y, com penalidade quantificada em Z'' é contribuição legítima e
publicável. Ela responde uma pergunta que a comunidade tem e que ninguém respondeu
com dados.

A equipe precisa ouvir isso do orientador \emph{antes} de começar. Equipes que não
ouvem tendem, no fim do semestre e sob pressão de prazo, a selecionar dados
favoráveis, estreitar faixas de análise ou apresentar resultados preliminares como
definitivos. A prevenção é declarar, no dia um, que o desfecho desfavorável tem valor.

%=========================================================================
\section{Produtos esperados}
%=========================================================================

\begin{table}[htbp]
\centering
\caption{Produtos esperados ao término do horizonte de 16 semanas.}
\label{tab:produtos}
\small
\begin{tabularx}{\textwidth}{@{}L{4.6cm} X@{}}
\toprule
\textbf{Produto} & \textbf{Descrição} \\
\midrule
Implementação em MAX 10 & Projeto Quartus completo, versionado, reprodutível, com restrições documentadas. \\
Relatório de caracterização da régua & Passo, dispersão, não-linearidade e deriva térmica da referência temporal adotada. \\
Conjunto de dados de campanha & Séries de \emph{slack}, temperatura e tensão, com proveniência documentada. \\
Análise comparativa & $\sigma_\text{resid}$ e degrau resolvível em ambas as plataformas, com incertezas. \\
Relatório técnico interno & Documento completo do port, incluindo decisões de projeto e desfechos de portão. \\
Submissão a conferência & Artigo sobre portabilidade entre fabricantes de sensores de \emph{slack} para SLM. \\
Material de formação & Trilhas de estudo consolidadas, reutilizáveis para as próximas turmas do laboratório. \\
\bottomrule
\end{tabularx}
\end{table}

\clearpage
%=========================================================================
\appendix
\section{Glossário}
%=========================================================================

\begin{longtable}{@{}L{3.0cm} L{12.1cm}@{}}
\toprule
\textbf{Termo} & \textbf{Definição} \\
\midrule
\endhead
AF & Fator de aceleração. Razão entre a taxa de degradação sob condição de estresse e sob condição de uso. \\
AVS & \emph{Adaptive Voltage Scaling}. Ajuste dinâmico de tensão de alimentação com base em medida de margem. \\
Auto-\emph{tuning} & Neste projeto, a capacidade do sensor de localizar e rastrear automaticamente a borda de falha do caminho sob teste. \\
CFM / UFM & Blocos de memória \emph{flash} de configuração e de usuário integrados ao \emph{die} do MAX 10. \\
DUT & \emph{Device Under Test}. Dispositivo sob ensaio. \\
$E_a$ & Energia de ativação do mecanismo de degradação, parâmetro da relação de Arrhenius. \\
FOPDT & \emph{First Order Plus Dead Time}. Modelo de planta de primeira ordem com atraso de transporte. \\
HCI & \emph{Hot Carrier Injection}. Mecanismo de degradação por portadores energéticos. \\
JTAG UART & Periférico Intel que transporta dados seriais pela cadeia JTAG, dispensando ponte USB-serial dedicada. \\
LE & \emph{Logic Element}. Célula lógica básica do MAX 10, com LUT de 4 entradas. \\
MMCM & \emph{Mixed-Mode Clock Manager}. Gerenciador de \emph{clock} da família Xilinx 7, com deslocamento de fase de passo $T_\text{VCO}/56$. \\
NBTI & \emph{Negative Bias Temperature Instability}. Mecanismo dominante de degradação em transistores PMOS. \\
PVT & Processo, tensão e temperatura. Conjunto de fontes de variação reversível. \\
RUL & \emph{Remaining Useful Life}. Vida útil remanescente estimada. \\
SIMC & \emph{Simple Internal Model Control}. Conjunto de regras de sintonia de PID. \\
\emph{Slack} & Margem temporal remanescente em um caminho síncrono; grandeza medida pelo sensor. \\
SLM & \emph{Silicon Lifecycle Management}. Gerenciamento de ciclo de vida do silício. \\
SSR & \emph{Solid State Relay}. Relé de estado sólido usado no acionamento do forno. \\
TDC & \emph{Time-to-Digital Converter}. Conversor tempo-digital. \\
TSD & \emph{Temperature Sensing Diode}. Diodo térmico integrado, lido pelo ADC do MAX 10. \\
Vernier & Técnica de medição temporal por diferença entre duas escalas próximas, com resolução efetiva igual à diferença. \\
WYSIWYG & Primitiva de baixo nível do fluxo Intel que instancia diretamente uma célula do dispositivo. \\
XADC & Bloco de conversão analógico-digital integrado da família Xilinx 7, com monitores internos de temperatura e alimentação. \\
\bottomrule
\end{longtable}

%=========================================================================
\section{Lista de verificação de aceite por frente}
%=========================================================================

Esta lista é usada nas reuniões de portão. Cada item é respondido com evidência
documental, não com relato verbal.

\campo{F1 --- Régua temporal e RTL}
\begin{itemize}
  \item[$\square$] Passo de fase medido em \emph{hardware}, com dispersão reportada.
  \item[$\square$] Configuração de VCO documentada e forçada, não deixada a critério da ferramenta.
  \item[$\square$] Diagrama de estados do laço portado, conferido contra o documento da trilha E3.
  \item[$\square$] Banco de testes cobrindo extremidade de faixa e reaquisição.
  \item[$\square$] Parecer sobre Vernier com dados de \emph{jitter} relativo.
  \item[$\square$] Laço rastreando borda deslocada por variação térmica, em \emph{hardware}.
\end{itemize}

\campo{F2 --- Implementação física}
\begin{itemize}
  \item[$\square$] \texttt{.qsf} integralmente comentado, com justificativa por atribuição.
  \item[$\square$] Otimizações de \emph{fitter} desabilitadas e listadas.
  \item[$\square$] Dez compilações com posicionamento idêntico, verificado automaticamente.
  \item[$\square$] Capturas do Chip Planner arquivadas.
  \item[$\square$] Versão do Quartus registrada no repositório.
\end{itemize}

\campo{F3 --- Instrumentação PVT}
\begin{itemize}
  \item[$\square$] Calibração do TSD em cinco pontos, com curva de erro residual.
  \item[$\square$] Esquema de intercalação TSD / canais externos documentado, com erro de sincronismo quantificado.
  \item[$\square$] Revisão de projeto da modificação de placa registrada antes da execução.
  \item[$\square$] Documentação fotográfica completa da intervenção.
  \item[$\square$] Leitura simultânea de temperatura e tensão chegando ao \emph{host}.
\end{itemize}

\campo{F4 --- Telemetria e aquisição}
\begin{itemize}
  \item[$\square$] Decisão de transporte registrada com justificativa.
  \item[$\square$] Compatibilidade de formato CSV verificada contra o \emph{pipeline} existente.
  \item[$\square$] Política de \emph{timestamp} documentada, com erro de correlação quantificado.
  \item[$\square$] Ensaio de \SI{72}{\hour} concluído, com relatório de perdas.
\end{itemize}

\campo{F5 --- Bancada}
\begin{itemize}
  \item[$\square$] Modelo FOPDT reidentificado para o novo DUT.
  \item[$\square$] PID ressintonizado e validado em campanha de \SI{12}{\hour}.
  \item[$\square$] Envelope de estresse documentado com os três limites justificados.
  \item[$\square$] Procedimento de verificação de integridade da configuração definido e testado.
  \item[$\square$] Procedimento operacional padrão de campanha escrito.
\end{itemize}

\campo{F6 --- Caracterização e análise}
\begin{itemize}
  \item[$\square$] Reprodução dos números publicados do Artix-7 confirmada.
  \item[$\square$] $\sigma_\text{resid}$ do MAX 10 reportado em picossegundos com incerteza.
  \item[$\square$] Degrau resolvível declarado para $k=3$ e ao menos um $k$ alternativo.
  \item[$\square$] Proveniência de cada número rastreável até o arquivo bruto.
  \item[$\square$] Tabela comparativa entre plataformas concluída.
\end{itemize}

%=========================================================================
\section{Bibliografia de estudo}
%=========================================================================

Lista de partida para a Fase 0. Não é exaustiva; cada trilha deve ampliá-la e registrar
as adições no repositório.

\campo{Normas e documentação de dispositivo}
\begin{itemize}
  \item JEDEC JESD22-A108 --- \emph{Temperature, Bias and Operating Life}.
  \item JEDEC JESD91 --- métodos de extração de fator de aceleração.
  \item Intel, \emph{MAX 10 Clocking and PLL User Guide}.
  \item Intel, \emph{MAX 10 Analog to Digital Converter User Guide}.
  \item Intel, \emph{MAX 10 FPGA Device Datasheet} --- características de ADC, sensor
        térmico, limites térmicos e retenção de memória de configuração.
  \item Intel, \emph{Quartus Prime Handbook} --- atribuições do \emph{fitter},
        restrições de posicionamento, Chip Planner.
  \item Xilinx, UG472 (\emph{7 Series FPGAs Clocking Resources}) e UG480 (\emph{XADC})
        --- para comparação com a plataforma de referência.
\end{itemize}

\campo{Trabalho anterior do grupo}
\begin{itemize}
  \item Artigo SBCCI 2025 --- validação do sensor de \emph{aging} auto-\emph{tuning}.
  \item Artigo SBCCI 2026 --- arquiteturas de controle térmico de burn-in e seu efeito
        sobre o fator de aceleração entregue e sobre a variância PVT do sensor.
  \item Trabalho final de TIP7266 --- caracterização informacional de envelhecimento e
        ruído em sensores de \emph{slack} embarcados para SLM.
  \item Código-fonte do sensor, do software de aquisição e do \emph{pipeline} de análise.
\end{itemize}

\campo{Literatura de referência}
\begin{itemize}
  \item Sensores in-situ de margem temporal: família Razor e derivados; canários de
        atraso; monitores por oscilador em anel.
  \item Conversores tempo-digital em FPGA: linha de atraso, Vernier, batimento de
        frequência.
  \item Gerenciamento de ciclo de vida do silício: artigos de revisão sobre SLM, AVS
        e estimação de vida útil remanescente.
  \item Confiabilidade de FPGAs comerciais sob envelhecimento acelerado.
\end{itemize}


\end{document}
